% Source: https://artofproblemsolving.com/community/c3h1411254_formula_for_cones
% Stable identity: aops-topic-1411254-q01
A right circular cone has base radius $r$ and slant height $\ell$. A regular $n$-gon is inscribed in its base circle, and its vertices are joined to the apex, forming a pyramid inscribed in the cone.
\begin{enumerate}
\item[(a)] Show that the lateral area of the pyramid is $A_n=nr\sin\dfrac{\pi}{n}\sqrt{\ell^2-r^2\sin^2\dfrac{\pi}{n}}$.
\item[(b)] Show that $A_n$ tends to $\pi r\ell$ as $n\to\infty$, and explain why one cannot simply substitute the perimeter $2\pi r$ into the pyramid formula ``lateral area $=\dfrac{1}{2}\times$ perimeter $\times$ slant height'' without saying which slant height is meant.
\end{enumerate}
